\documentclass[a4paper]{article}

%------------------------------------------------
% الحزم
%------------------------------------------------

\usepackage{enumitem}
\usepackage[margin=2.5cm]{geometry}
\usepackage{amsmath,amsfonts,amssymb}
\usepackage[charter]{mathdesign}
\usepackage{graphicx}
\usepackage{array}
\usepackage{xcolor}
\usepackage{needspace}
\usepackage{fancyhdr}
\usepackage{polyglossia}

%------------------------------------------------
% اللغات والخطوط
%------------------------------------------------

\setdefaultlanguage[
    calendar=gregorian,
    numerals=maghrib
]{arabic}

\setotherlanguage{english}

% إذا لم تكن الخطوط مثبتة في الجهاز، تُؤخذ من المجلد ../../fonts/
\IfFontExistsTF{Amiri}{%
    \newfontfamily\arabicfont[Script=Arabic, Scale=1.2]{Amiri}
}{%
    \newfontfamily\arabicfont[
        Script=Arabic, Scale=1.2,
        Path=../../fonts/, Extension=.ttf,
        UprightFont=*-Regular, BoldFont=*-Bold,
        ItalicFont=*-Italic, BoldItalicFont=*-BoldItalic
    ]{Amiri}
}

\IfFontExistsTF{AlBattar}{%
    \newfontfamily\arabicfontsf[Script=Arabic, Scale=1.3]{AlBattar}
    \newfontfamily\khatt[Script=Arabic, Scale=1.3]{AlBattar}
}{%
    \newfontfamily\arabicfontsf[
        Script=Arabic, Scale=1.3, Path=../../fonts/,
        BoldFont=ae_AlBattar.ttf
    ]{ae_AlBattar.ttf}
    \newfontfamily\khatt[
        Script=Arabic, Scale=1.3, Path=../../fonts/,
        BoldFont=ae_AlBattar.ttf
    ]{ae_AlBattar.ttf}
}

%------------------------------------------------
% نسخة الطالب / نسخة التصحيح
% لإنتاج نسخة التصحيح: عرّف \withanswers قبل \input لهذا الملف
%------------------------------------------------

\newif\ifcorrection
\ifdefined\withanswers \correctiontrue \else \correctionfalse \fi

\newcommand{\emptybox}{$\square$}
\ifcorrection
    \newcommand{\correctbox}{\textcolor{green!50!black}{$\boxtimes$}}
    \newcommand{\correct}[1]{\textcolor{green!40!black}{\textbf{#1}}}
    \newcommand{\explanation}[1]{%
        \par\vspace{3pt}
        \fcolorbox{gray!60}{gray!8}{\parbox{\dimexpr\linewidth-2\fboxsep-2\fboxrule}{%
        \small\textbf{التعليل:} #1}}}
\else
    \newcommand{\correctbox}{$\square$}
    \newcommand{\correct}[1]{#1}
    \newcommand{\explanation}[1]{}
\fi

\newcommand{\pts}[1]{\textbf{(#1 ن)}}

%------------------------------------------------
% إعداد الصفحة
%------------------------------------------------

\renewcommand{\footrulewidth}{1pt}
\renewcommand{\headrulewidth}{1pt}

\rhead{\khatt المدرسة العليا للأساتذة ببوزريعة}
\lhead{\khatt ملحقة بومرداس}

\cfoot{\thepage}
\lfoot{\khatt الأستاذ: د. بن عبد الرزاق يوسف}

\parindent 0pt
\setlength{\headsep}{15pt}
\arraycolsep=1.4pt

\pagestyle{fancy}

\setlist[enumerate,1]{label=\textbf{السؤال \arabic*:}, leftmargin=*, itemsep=12pt, align=left}

%================================================
% بداية الوثيقة
%================================================

\begin{document}

\begin{center}

{\Large\khatt\textbf{الواجب الرابع}\ifcorrection\ \textbf{-- التصحيح النموذجي}\fi}

\vspace{0.2cm}

{\Large\khatt\textbf{في الذكاء الاصطناعي}}

\end{center}

\rule{\textwidth}{1pt}

\ifcorrection\else
\textbf{الاسم واللقب:} \dotfill \quad \textbf{الفوج:} \makebox[3cm]{\dotfill}

\vspace{0.3cm}
\fi

\begin{center}

{\khatt\textbf{أسئلة متعددة الاختيارات (20 سؤالا -- 15 نقطة)}}

\end{center}

\rule{\textwidth}{1pt}

\textbf{التعليمة:} ضع علامة \textbf{$\times$} في المربع المقابل للإجابة الصحيحة. لكل سؤال إجابة صحيحة واحدة فقط.

\vspace{0.3cm}

\begin{enumerate}

\needspace{9\baselineskip}\item
\textbf{في نقاش حول الذكاء الاصطناعي في التعليم، قال طالب: «الأداة هي المشكلة، فهي التي تدفع الطلبة إلى الغش». أي ردّ يوافق ما جاء في الدرس؟} \hfill \pts{0,75}
\begin{itemize}[leftmargin=1.6cm, labelsep=0.3cm, itemsep=2pt, topsep=4pt]
    \item[\emptybox\ \textbf{أ)}] الغش مسؤولية الشركات المطوّرة وحدها، فالطالب يستعمل ما يُتاح له، ولا يُحاسَب على طريقة استعماله لأي أداة.
    \item[\correctbox\ \textbf{ب)}] \correct{الأداة لا تقرر وحدها كيف تُستعمل؛ الإنسان هو الذي يختار أن يستعملها بأمانة أو للغش، ولهذا نتحدث عن الأخلاقيات.}
    \item[\emptybox\ \textbf{ج)}] لا علاقة للأخلاقيات بالذكاء الاصطناعي، لأن الأدوات التقنية محايدة دائمًا ولا يترتب على استعمالها أي مسؤولية.
    \item[\emptybox\ \textbf{د)}] الأداة مسؤولة فعلًا عن الغش، لأنها مصمّمة لتعطي إجابات كاملة، ولذلك يجب منعها نهائيًا من كل المؤسسات التعليمية.
\end{itemize}
\explanation{يؤكد الدرس أن الأداة قوية وسريعة لكنها لا تقرر كيف نستعملها. الإنسان هو الذي يختار، ومن هنا تأتي أهمية الأخلاقيات.}

%------------------------------------------------

\needspace{9\baselineskip}\item
\textbf{قدّم أستاذ متربص في القسم معلومة تاريخية خاطئة أخذها من أداة ذكاء اصطناعي دون تحقق. من المسؤول عن هذا الخطأ حسب الدرس؟} \hfill \pts{0,75}
\begin{itemize}[leftmargin=1.6cm, labelsep=0.3cm, itemsep=2pt, topsep=4pt]
    \item[\emptybox\ \textbf{أ)}] الشركة المطوّرة للأداة، لأنها لم تمنع ظهور المعلومات الخاطئة، فالمسؤولية تقع كلها على من صنع الأداة.
    \item[\emptybox\ \textbf{ب)}] التلاميذ، لأنهم لم ينتبهوا إلى الخطأ أثناء الحصة، وكان عليهم أن يتحققوا بأنفسهم من كل ما يقدّمه الأستاذ.
    \item[\emptybox\ \textbf{ج)}] الأداة وحدها، لأنها هي التي ولّدت المعلومة الخاطئة، والأستاذ نقلها بحسن نية دون أن يتدخل في صياغتها.
    \item[\correctbox\ \textbf{د)}] \correct{الأستاذ نفسه، لأن الاستعانة بالأداة لا تنقل المسؤولية إليها، والمستخدم مسؤول عن كل معلومة يعتمد عليها.}
\end{itemize}
\explanation{مبدأ المسؤولية: المستخدم مسؤول عن استعمال الأداة وعن كل معلومة يعتمد عليها. إذا قدّمت معلومة خاطئة، فأنت المسؤول لا الأداة.}

%------------------------------------------------

\needspace{9\baselineskip}\item
\textbf{سأل أستاذ طالبته: «كيف استعملتِ الذكاء الاصطناعي في هذا البحث بالضبط؟»، فشرحت أنها استعملته لاقتراح خطة أولية، ثم بحثت في المصادر وكتبت الفصول بنفسها. أي مبدأ يتجسّد في هذا الشرح المفصّل لطريقة الاستعمال؟} \hfill \pts{0,75}
\begin{itemize}[leftmargin=1.6cm, labelsep=0.3cm, itemsep=2pt, topsep=4pt]
    \item[\emptybox\ \textbf{أ)}] المسؤولية، لأنها تنقل عبء الأخطاء من الطالبة إلى الأداة.
    \item[\emptybox\ \textbf{ب)}] الانتحال، لأنها اعترفت باستعمال أداة في إنجاز البحث.
    \item[\correctbox\ \textbf{ج)}] \correct{الشفافية، لأنها توضّح طريقة استعمال الأداة عند الحاجة.}
    \item[\emptybox\ \textbf{د)}] التلخيص، لأنها أعادت كتابة أفكارها بكلمات أقل من الأصل.
\end{itemize}
\explanation{الشفافية هي توضيح طريقة استعمال الذكاء الاصطناعي عند الحاجة، أي ماذا فعلت الأداة وماذا فعل الطالب.}

%------------------------------------------------

\needspace{9\baselineskip}\item
\textbf{ما التعريف الصحيح لمبدأ «الإفصاح» كما ورد في الدرس؟} \hfill \pts{0,75}
\begin{itemize}[leftmargin=1.6cm, labelsep=0.3cm, itemsep=2pt, topsep=4pt]
    \item[\emptybox\ \textbf{أ)}] ذكر اسم الأستاذ المشرف في الصفحة الأولى من البحث بدل اسم الطالب.
    \item[\emptybox\ \textbf{ب)}] نشر البحث على الإنترنت حتى يطّلع عليه أكبر عدد ممكن من القرّاء.
    \item[\correctbox\ \textbf{ج)}] \correct{التصريح باستعمال أداة الذكاء الاصطناعي عندما يكون ذلك مطلوبًا أو مناسبًا.}
    \item[\emptybox\ \textbf{د)}] إخفاء استعمال الأداة إذا كانت الصياغة النهائية جيدة ولا يلاحظ أحد الفرق.
\end{itemize}
\explanation{الإفصاح هو التصريح باستعمال أداة الذكاء الاصطناعي عندما يكون ذلك مطلوبًا أو مناسبًا، ويكون غالبًا بجملة إفصاح قصيرة وصادقة.}

%------------------------------------------------

\needspace{9\baselineskip}\item
\textbf{أي تعريف يعبّر بدقة عن «الانتحال العلمي»؟} \hfill \pts{0,75}
\begin{itemize}[leftmargin=1.6cm, labelsep=0.3cm, itemsep=2pt, topsep=4pt]
    \item[\emptybox\ \textbf{أ)}] أن يكتب الشخص بحثه بأسلوبه الخاص بعد قراءة عدة مصادر مختلفة.
    \item[\correctbox\ \textbf{ب)}] \correct{أن ينسب الشخص إلى نفسه عملًا أو أفكارًا ليست له، دون ذكر مصدرها.}
    \item[\emptybox\ \textbf{ج)}] أن يستعمل الشخص أداة ذكاء اصطناعي لفهم مفهوم صعب قبل الكتابة.
    \item[\emptybox\ \textbf{د)}] أن يقتبس الشخص فقرة من كتاب مع ذكر المؤلف والصفحة بشكل صحيح.
\end{itemize}
\explanation{الانتحال العلمي هو نسبة عمل الغير أو أفكاره إلى النفس دون ذكر المصدر. الاقتباس مع ذكر المصدر ليس انتحالًا.}

%------------------------------------------------

\needspace{9\baselineskip}\item
\textbf{نسخ طالب بحثًا كاملًا أنتجته أداة ذكاء اصطناعي، ثم كتب اسمه عليه وقدّمه للأستاذ. كيف يصنّف الدرس هذا السلوك؟} \hfill \pts{0,75}
\begin{itemize}[leftmargin=1.6cm, labelsep=0.3cm, itemsep=2pt, topsep=4pt]
    \item[\emptybox\ \textbf{أ)}] استعمال مقبول، لأن الطالب هو الذي كتب الأمر، فالنص الناتج يُعدّ ملكًا له بالكامل.
    \item[\emptybox\ \textbf{ب)}] سلوك محايد، لأن القوانين لا تتحدث عن الذكاء الاصطناعي، فلا يمكن وصفه بالانتحال.
    \item[\correctbox\ \textbf{ج)}] \correct{انتحال بالذكاء الاصطناعي: العمل ليس عمله، ولم يتحقق من صحته، ولم يصرّح بمصدره.}
    \item[\emptybox\ \textbf{د)}] استعمال مقبول بشروط، يكفي فيه أن يغيّر الطالب الخط وحجمه قبل تقديم البحث.
\end{itemize}
\explanation{هذا مثال الدرس على الانتحال بالذكاء الاصطناعي: تقديم نص أنتجته الأداة على أنه عمل شخصي، دون تحقق ودون تصريح.}

%------------------------------------------------

\needspace{9\baselineskip}\item
\textbf{استعملت طالبة أداة ذكاء اصطناعي لتفهم قاعدة نحوية صعبة في الممنوع من الصرف، ثم أنجزت تمارينها بنفسها. كيف يُصنّف هذا الاستعمال؟} \hfill \pts{0,75}
\begin{itemize}[leftmargin=1.6cm, labelsep=0.3cm, itemsep=2pt, topsep=4pt]
    \item[\correctbox\ \textbf{أ)}] \correct{مقبول، لأنها استعملت الأداة للفهم، والعمل المقدَّم من إنجازها.}
    \item[\emptybox\ \textbf{ب)}] مقبول بشروط، بشرط أن تنسخ شرح الأداة في كراسها حرفيًا.
    \item[\emptybox\ \textbf{ج)}] غير مسؤول، لأنها لم تكتب جملة إفصاح في كراس التمارين.
    \item[\emptybox\ \textbf{د)}] مرفوض، لأن أي استعمال للأداة في الدراسة يُعدّ انتحالًا.
\end{itemize}
\explanation{في مقياس الاستعمال: استعمال الأداة لفهم مفهوم أو قاعدة صعبة استعمال مقبول.}

%------------------------------------------------

\needspace{9\baselineskip}\item
\textbf{طلب طالب من أداة أفكارًا لموضوع إنشاء، ثم اختار منها ما يناسبه وكتب الموضوع كله بأسلوبه. ما الموقف الأخلاقي من هذا الاستعمال؟} \hfill \pts{0,75}
\begin{itemize}[leftmargin=1.6cm, labelsep=0.3cm, itemsep=2pt, topsep=4pt]
    \item[\emptybox\ \textbf{أ)}] انتحال مرفوض، لأن الأفكار الأولى جاءت من الأداة لا من الطالب.
    \item[\emptybox\ \textbf{ب)}] غير مسؤول، لأن الطالب لم يقبل كل الأفكار التي اقترحتها الأداة.
    \item[\emptybox\ \textbf{ج)}] مرفوض، لأن طلب الأفكار أخطر من نسخ إجابة كاملة من الأداة.
    \item[\correctbox\ \textbf{د)}] \correct{مقبول، لأن الأداة اقترحت أفكارًا فقط، والكتابة من عمل الطالب.}
\end{itemize}
\explanation{في جدول الدرس: طلب أفكار من الأداة ثم كتابة الإجابة بنفسي استعمال مقبول.}

%------------------------------------------------

\needspace{9\baselineskip}\item
\textbf{أنشأ طالب مسودة لبحثه بأداة ذكاء اصطناعي، ثم راجعها وتحقق من معلوماتها في المصادر، وأضاف جملة إفصاح. كيف يصنّف الدرس هذا الاستعمال؟} \hfill \pts{0,75}
\begin{itemize}[leftmargin=1.6cm, labelsep=0.3cm, itemsep=2pt, topsep=4pt]
    \item[\correctbox\ \textbf{أ)}] \correct{مقبول بشروط: المراجعة والتحقق والإفصاح هي التي تجعله مقبولًا.}
    \item[\emptybox\ \textbf{ب)}] مقبول دون أي شرط: ولا حاجة إلى المراجعة أو التحقق أو الإفصاح.
    \item[\emptybox\ \textbf{ج)}] مرفوض دائمًا: لأن إنشاء المسودة بالأداة انتحال مهما فعل بعد ذلك.
    \item[\emptybox\ \textbf{د)}] غير مسؤول: لأن جملة الإفصاح تكشف ضعف الطالب أمام أستاذه.
\end{itemize}
\explanation{في الدرس: إنشاء مسودة بالأداة ثم مراجعتها والتحقق منها والإفصاح استعمال «مقبول بشروط». غياب هذه الشروط يحوّله إلى استعمال مرفوض.}

%------------------------------------------------

\needspace{9\baselineskip}\item
\textbf{اعتمد طالب على تواريخ وأسماء قدّمتها أداة ذكاء اصطناعي، وأدرجها في بحثه دون أن يتحقق منها في أي مصدر. كيف يصف الدرس هذا السلوك؟} \hfill \pts{0,75}
\begin{itemize}[leftmargin=1.6cm, labelsep=0.3cm, itemsep=2pt, topsep=4pt]
    \item[\correctbox\ \textbf{أ)}] \correct{غير مسؤول، لأنه اعتمد على معلومات الأداة دون التحقق منها.}
    \item[\emptybox\ \textbf{ب)}] شفاف، لأن الطالب استعمل الأداة بطريقة واضحة ومباشرة.
    \item[\emptybox\ \textbf{ج)}] مقبول بشروط، يكفي فيه أن تكون المعلومات مكتوبة بثقة.
    \item[\emptybox\ \textbf{د)}] مقبول، لأن التواريخ والأسماء لا تحتاج إلى أي تحقق إضافي.
\end{itemize}
\explanation{الاعتماد على معلومات الأداة دون التحقق منها مصنّف في الدرس على أنه «غير مسؤول»، خاصة أن الأداة قد تخطئ أو تهلوس.}

%------------------------------------------------

\needspace{9\baselineskip}\item
\textbf{ما التصنيف الذي يعطيه الدرس لنسخ إجابة كاملة من أداة ذكاء اصطناعي وتقديمها باسم الطالب؟} \hfill \pts{0,75}
\begin{itemize}[leftmargin=1.6cm, labelsep=0.3cm, itemsep=2pt, topsep=4pt]
    \item[\emptybox\ \textbf{أ)}] استعمال مقبول.
    \item[\emptybox\ \textbf{ب)}] استعمال شفاف.
    \item[\correctbox\ \textbf{ج)}] \correct{انتحال مرفوض.}
    \item[\emptybox\ \textbf{د)}] مقبول بشروط.
\end{itemize}
\explanation{نسخ إجابة كاملة وتقديمها باسمي يقع في أقصى مقياس الاستعمال: انتحال مرفوض.}

%------------------------------------------------

\needspace{9\baselineskip}\item
\textbf{طُلب من طلبة كتابة جملة إفصاح في آخر بحثهم. أي جملة تحقق شروط الإفصاح الجيد؟} \hfill \pts{0,75}
\begin{itemize}[leftmargin=1.6cm, labelsep=0.3cm, itemsep=2pt, topsep=4pt]
    \item[\correctbox\ \textbf{أ)}] \correct{استعملتُ أداة \textenglish{ChatGPT} لتنظيم أفكار الفصل الثاني واقتراح خطة أولية، وراجعتُ كل المعلومات وتحققت منها في المصادر المذكورة، والصياغة النهائية من عملي.}
    \item[\emptybox\ \textbf{ب)}] لم أستعمل أي أداة ذكاء اصطناعي في هذا البحث، مع أنني استعملت أداة لكتابة المقدمة والخاتمة، لأن المقدمة والخاتمة لا تُعدّان جزءًا من البحث.
    \item[\emptybox\ \textbf{ج)}] هذا البحث كتبه الذكاء الاصطناعي بالكامل، وأنا غير مسؤول عن أي خطأ يرد فيه، فعلى القارئ أن يتحقق بنفسه من كل معلومة قبل أن يعتمد عليها.
    \item[\emptybox\ \textbf{د)}] قد أكون استعملت بعض الأدوات الحديثة في مكان ما من هذا البحث، لكن ذلك غير مهم ولا يؤثر في قيمة العمل، لذلك لا داعي لذكر التفاصيل هنا.
\end{itemize}
\explanation{جملة الإفصاح الجيدة قصيرة وصادقة: تذكر الأداة، وما استُعملت فيه، وتؤكد التحقق، وتوضح أن الصياغة النهائية من عمل الطالب. هذا هو نموذج الدرس.}

%------------------------------------------------

\needspace{9\baselineskip}\item
\textbf{أين توضع جملة الإفصاح عادة، وما صفاتها حسب الدرس؟} \hfill \pts{0,75}
\begin{itemize}[leftmargin=1.6cm, labelsep=0.3cm, itemsep=2pt, topsep=4pt]
    \item[\emptybox\ \textbf{أ)}] في قائمة المراجع، مكان اسم المؤلف الأول فيها.
    \item[\correctbox\ \textbf{ب)}] \correct{في آخر البحث أو في مقدمته، وتكون قصيرة وصادقة.}
    \item[\emptybox\ \textbf{ج)}] في رسالة خاصة إلى الأستاذ، ولا تُكتب في البحث.
    \item[\emptybox\ \textbf{د)}] في وسط الفصل الأول، وتكون طويلة ومفصّلة جدًا.
\end{itemize}
\explanation{حسب الدرس: جملة الإفصاح قصيرة وصادقة، نضعها في آخر البحث أو في مقدمته.}

%------------------------------------------------

\needspace{9\baselineskip}\item
\textbf{أي عبارة من العبارات الآتية تُعدّ من «القواعد الذهبية للاستعمال الأخلاقي» التي وردت في الدرس؟} \hfill \pts{0,75}
\begin{itemize}[leftmargin=1.6cm, labelsep=0.3cm, itemsep=2pt, topsep=4pt]
    \item[\emptybox\ \textbf{أ)}] أقبل كل معلومة تقدّمها الأداة دون تحقق.
    \item[\correctbox\ \textbf{ب)}] \correct{أكتب العمل النهائي بأسلوبي وأفكاري.}
    \item[\emptybox\ \textbf{ج)}] أخفي استعمال الأداة إذا لم يسألني أحد.
    \item[\emptybox\ \textbf{د)}] أنسخ الإجابة إذا كانت صياغتها أفضل من صياغتي.
\end{itemize}
\explanation{القواعد الذهبية في الدرس: أتحقق من كل معلومة مهمة، أكتب العمل النهائي بأسلوبي وأفكاري، أصرّح باستعمال الأداة عند الحاجة، أحترم قوانين المؤسسة وتعليمات الأستاذ.}

%------------------------------------------------

\needspace{9\baselineskip}\item
\textbf{اقترح الدرس قاعدة بسيطة يعلّمها أستاذ المستقبل لتلاميذه حول استعمال الذكاء الاصطناعي. ما هي؟} \hfill \pts{0,75}
\begin{itemize}[leftmargin=1.6cm, labelsep=0.3cm, itemsep=2pt, topsep=4pt]
    \item[\emptybox\ \textbf{أ)}] «الأداة لا تخطئ أبدًا، فأنقل منها دون تردد».
    \item[\emptybox\ \textbf{ب)}] «الأداة تفكر بدلًا مني، فأرتاح من عناء التعلّم».
    \item[\correctbox\ \textbf{ج)}] \correct{«الأداة تساعدني على التفكير، ولا تفكر بدلًا مني».}
    \item[\emptybox\ \textbf{د)}] «الأداة ممنوعة دائمًا، فلا أقترب منها مطلقًا».
\end{itemize}
\explanation{هذه القاعدة من الدرس تلخّص الاستعمال الأخلاقي: الأداة وسيلة مساعدة على التفكير، لا بديل عنه.}

%------------------------------------------------

\needspace{9\baselineskip}\item
\textbf{منع أستاذ استعمال أدوات الذكاء الاصطناعي في واجب معيّن. استعملها طالب رغم ذلك، ثم قال: «لا مشكلة، فقد كتبت جملة إفصاح». ما الموقف الصحيح؟} \hfill \pts{0,75}
\begin{itemize}[leftmargin=1.6cm, labelsep=0.3cm, itemsep=2pt, topsep=4pt]
    \item[\emptybox\ \textbf{أ)}] الطالب محقّ تمامًا، لأن جملة الإفصاح تجعل أي استعمال مقبولًا مهما كانت تعليمات الأستاذ.
    \item[\emptybox\ \textbf{ب)}] لا موقف أخلاقي هنا، لأن تعليمات الأستاذ تخص الواجبات الورقية فقط، ولا تشمل الأدوات الرقمية.
    \item[\emptybox\ \textbf{ج)}] الطالب محقّ إذا كانت جملة الإفصاح طويلة، لأن طولها يعوّض مخالفة التعليمات التي حددها الأستاذ.
    \item[\correctbox\ \textbf{د)}] \correct{الإفصاح لا يبرّر مخالفة التعليمات؛ من القواعد الذهبية احترام قوانين المؤسسة وتعليمات الأستاذ.}
\end{itemize}
\explanation{الإفصاح مبدأ مهم، لكنه لا يلغي القواعد الأخرى. احترام قوانين المؤسسة وتعليمات الأستاذ من القواعد الذهبية في الدرس.}

%------------------------------------------------

\needspace{9\baselineskip}\item
\textbf{اكتُشف خطأ علمي في بحث طالب، فقال: «لست مسؤولًا، فالأداة هي التي كتبت هذه الفقرة». بماذا يُردّ عليه حسب الدرس؟} \hfill \pts{0,75}
\begin{itemize}[leftmargin=1.6cm, labelsep=0.3cm, itemsep=2pt, topsep=4pt]
    \item[\correctbox\ \textbf{أ)}] \correct{الاستعانة بالذكاء الاصطناعي لا تنقل المسؤولية إلى الأداة؛ فمن يقدّم العمل مسؤول عن كل معلومة فيه.}
    \item[\emptybox\ \textbf{ب)}] هو محقّ، لأن المسؤولية تنتقل تلقائيًا إلى الأداة متى كانت هي التي كتبت الفقرة التي تحتوي على الخطأ.
    \item[\emptybox\ \textbf{ج)}] هو محقّ إذا كان قد كتب جملة إفصاح، لأن الإفصاح ينقل المسؤولية كاملة من الطالب إلى الأداة المستعملة.
    \item[\emptybox\ \textbf{د)}] المسؤولية مشتركة بالتساوي بين الطالب والأداة، فيُخصم نصف العلامة فقط بدل خصمها كاملة من الطالب.
\end{itemize}
\explanation{يؤكد الدرس: الاستعانة بالذكاء الاصطناعي لا تنقل المسؤولية إلى الأداة. المستخدم مسؤول عن كل معلومة يعتمد عليها.}

%------------------------------------------------

\needspace{9\baselineskip}\item
\textbf{نسخ طالب نصًا أنتجته أداة، ثم استعمل برنامجًا يغيّر بعض الكلمات بمرادفاتها، وقدّمه على أنه عمله دون تصريح. كيف يُقيَّم هذا السلوك؟} \hfill \pts{0,75}
\begin{itemize}[leftmargin=1.6cm, labelsep=0.3cm, itemsep=2pt, topsep=4pt]
    \item[\emptybox\ \textbf{أ)}] لم يعد انتحالًا، لأن تغيير الكلمات بمرادفاتها يجعل النص الجديد عملًا شخصيًا أصيلًا لا علاقة له بالأداة.
    \item[\emptybox\ \textbf{ب)}] سلوك شفاف، لأن الطالب استعمل برنامجين بدل برنامج واحد، والشفافية تعني تعدد الأدوات المستعملة في العمل.
    \item[\emptybox\ \textbf{ج)}] استعمال مقبول بشروط، لأن الطالب بذل جهدًا في استعمال برنامج المرادفات، وهذا يكفي ليُنسب العمل إليه.
    \item[\correctbox\ \textbf{د)}] \correct{يبقى انتحالًا، لأن الأفكار والعمل ليسا من إنجازه، وتغيير الكلمات لا يغيّر حقيقة نسبة عمل الغير إلى نفسه.}
\end{itemize}
\explanation{الانتحال بالذكاء الاصطناعي هو تقديم ما أنتجته الأداة على أنه عمل شخصي. تغيير الكلمات لا يجعل الأفكار أفكاره، ولا يعوّض غياب التصريح.}

%------------------------------------------------

\needspace{9\baselineskip}\item
\textbf{لخّص طالب درس الأخلاقيات في عبارة واحدة تصف الاستعمال الجيد للذكاء الاصطناعي. أي عبارة هي الأدق؟} \hfill \pts{0,75}
\begin{itemize}[leftmargin=1.6cm, labelsep=0.3cm, itemsep=2pt, topsep=4pt]
    \item[\emptybox\ \textbf{أ)}] أسأل، أقبل، أنقل، ولا أذكر الأداة.
    \item[\correctbox\ \textbf{ب)}] \correct{أفهم، أتحقق، أكتب بنفسي، وأصرّح.}
    \item[\emptybox\ \textbf{ج)}] أمتنع عن كل أداة في كل الأحوال.
    \item[\emptybox\ \textbf{د)}] أنسخ، أنسّق، أقدّم، وأنتظر العلامة.
\end{itemize}
\explanation{خلاصة الدرس: الاستعمال الجيد هو أن أفهم، وأتحقق، وأكتب بنفسي، وأصرّح باستعمال الأداة.}

%------------------------------------------------

\needspace{9\baselineskip}\item
\textbf{قال أحد الطلبة: «طريقة استعمالي للذكاء الاصطناعي شأن شخصي، ولا علاقة لها بمهنتي المستقبلية كأستاذ». لماذا يرفض الدرس هذا الرأي؟} \hfill \pts{0,75}
\begin{itemize}[leftmargin=1.6cm, labelsep=0.3cm, itemsep=2pt, topsep=4pt]
    \item[\emptybox\ \textbf{أ)}] لأن الأساتذة ممنوعون قانونيًا من استعمال أدوات الذكاء الاصطناعي في حياتهم الشخصية خارج المؤسسة.
    \item[\emptybox\ \textbf{ب)}] لأن الأداة تقيّم الأستاذ تلقائيًا وترسل تقريرًا عن طريقة استعماله إلى إدارة المؤسسة كل شهر.
    \item[\emptybox\ \textbf{ج)}] لأن التلاميذ يستطيعون الاطلاع على سجلّ محادثات أستاذهم مع الأداة، فيجب أن تكون كلها رسمية.
    \item[\correctbox\ \textbf{د)}] \correct{لأن أستاذ المستقبل قدوة لتلاميذه، فسلوكه مع هذه الأدوات درس في حد ذاته يتعلّمون منه الأمانة أو الغش.}
\end{itemize}
\explanation{يذكّر الدرس بأن أستاذ المستقبل قدوة لتلاميذه، فسلوكه مع أدوات الذكاء الاصطناعي درس في حد ذاته.}

\end{enumerate}

\vfill
\begin{center}
{\khatt بالتوفيق}
\end{center}

\end{document}
